% FILE: 06_contribution_to_thesis.tex
% PROJECT: otel-weaver-exp-001-custom-pi-registry
% THESIS ROLE: otel-semconv-evidence-surface-experiment

\section{Contribution to the Dissertation}
\label{sec:otel-weaver-exp-001-custom-pi-registry-contribution}

\subsection{Contribution to Prediction vs.~Coordination}
\label{sec:otel-weaver-exp-001-custom-pi-registry-contribution-prediction}

The dissertation's central argument distinguishes prediction from coordination: prediction models
future states from historical data; coordination enforces lawful state transitions in running
processes. The process intelligence architecture that the dissertation proposes is coordination-first:
it does not predict what processes will do; it determines whether what processes did was lawful,
and it does so from evidence that can be independently verified and replayed.

EXP-001 contributes to this argument by establishing that the coordination claim begins at the
feedstock layer. A process intelligence system that ingests unstructured telemetry cannot make
coordination claims: it can only make statistical claims about span distributions. By requiring
that admissible spans carry \texttt{process.pi.lifecycle} (one of six lawful lifecycle values),
\texttt{process.pi.activity.name} (a Petri net transition label), and
\texttt{process.pi.witness.hash} (a BLAKE3-sealed witness), EXP-001 ensures that the
coordination claim is grounded in evidence that was structurally enforced at emission time, not
inferred post hoc from unstructured logs.

The schema boundary is therefore the first coordination gate in the architecture. It precedes the
admission gate (\texttt{validate\_and\_project}), the bridging layer (\texttt{BridgeRx}), and
the conformance court (\texttt{wasm4pm}). All of those gates assume that their input has already
passed the feedstock schema check. The dissertation's coordination claim is only as strong as its
weakest feedstock constraint, and EXP-001's required-attribute set is what makes those constraints
machine-verifiable.

\subsection{Contribution to Process Evidence}
\label{sec:otel-weaver-exp-001-custom-pi-registry-contribution-evidence}

The dissertation argues that process evidence requires four properties: \emph{case identity}
(the event can be grouped with other events belonging to the same process instance),
\emph{activity identity} (the event maps to a named transition in a declared process model),
\emph{temporal ordering} (the event carries a timestamp that establishes its position in the
instance's lifecycle), and \emph{witness integrity} (the event's payload has been sealed by an
identifiable, verifiable witness). Standard OTel spans satisfy none of these properties by
default.

EXP-001 contributes a concrete instantiation of all four properties in a standard-conforming
schema:
\begin{itemize}
  \item Case identity $\to$ \texttt{process.pi.instance\_id} (required)
  \item Activity identity $\to$ \texttt{process.pi.activity.name} (required, Petri net transition label)
  \item Lifecycle phase $\to$ \texttt{process.pi.lifecycle} (required, six-value enumeration)
  \item Witness integrity $\to$ \texttt{process.pi.witness.hash} + \texttt{process.pi.witness.id} (both required)
\end{itemize}

The two recommended attributes (\texttt{process.pi.token.state\_before}/\texttt{state\_after})
add Petri net state evidence beyond the four foundational properties. This contribution is
empirical: the schema is a concrete artifact, not a design proposal, and it is validated against
the OTel Weaver toolchain rather than against the author's intent alone (subject to the PARTIAL
qualification that the \texttt{weaver registry validate} run has not been confirmed).

\subsection{Contribution to Manufacturing Architecture}
\label{sec:otel-weaver-exp-001-custom-pi-registry-contribution-manufacturing}

In the dissertation's manufacturing architecture, every artifact has a declared provenance chain:
a ggen template, a manifest, a pipeline execution, and a receipt. EXP-001 is the entry-layer
artifact in the OTel Weaver integration suite's manufacturing chain, and it demonstrates that a
schema definition --- typically treated as a configuration file with no formal provenance --- can
be placed under the same manufacturing discipline as code artifacts.

The ggen template (\texttt{pi-weaver-registry.yaml.ggen}) is the manufacturing input. The
manifest (\texttt{otel-weaver-source.manifest.yaml}) governs the pipeline. The checkpoint
sequence (PARTIAL $\to$ ALIVE $\to$ RUNTIME) records the pipeline's gate progression. The
declared BLAKE3 receipt chain (currently PARTIAL) is the provenance seal. This pattern ---
template $\to$ manifest $\to$ gates $\to$ receipt --- is the standard manufacturing discipline
applied to every artifact in the dissertation's evidence corpus. EXP-001's contribution is to
demonstrate that the discipline extends to schema artifacts as well as code artifacts.

\subsection{Contribution to Capital-Flow Infrastructure}
\label{sec:otel-weaver-exp-001-custom-pi-registry-contribution-capital-flow}

The dissertation's capital-flow infrastructure layer requires that process events in financial
and operational workflows be traceable, auditable, and replayable (AUTHOR THESIS). In the
capital-flow context, a process event is not merely a log entry: it is a claim about value
exchange, obligation fulfillment, or settlement state. Such claims require exactly the evidence
properties that EXP-001's schema enforces: case identity (which transaction instance?), activity
identity (which settlement step?), lifecycle phase (was this an initiation, a completion, or an
abort?), and witness integrity (who sealed this claim, and can the seal be verified?).

EXP-001 contributes the feedstock schema that would allow capital-flow process events to enter
the process intelligence conformance court as lawful evidence. Without this schema, capital-flow
telemetry from OTel-instrumented services is indistinguishable from operational telemetry at the
feedstock layer. With it, a capital-flow service can emit spans that carry the required attributes,
pass the admission gate, be projected to OCEL, and be replayed against a declared settlement
process model with measurable fitness and precision. This is the path from operational telemetry
to auditable capital-flow process intelligence, and EXP-001's schema is its entry point.
